# Tenchi o Kurau dialogue bank 11 -- kanji text, revision 14
#
# Generated by tools/kanji_layout.py from out/reconv/out/batch_*.txt.
# Edit the batch files (or this file) and rebuild.

b11_msg_000:
    "「ここまで来るとは　さすが<NL>劉備軍　まずは褒めておこう<PAGE>しかし　その程度の兵力で<NL>我らを倒すことができるかな？<END>"

b11_msg_001:
    "「我らの軍が敗れたのは<NL>初めてです<PAGE>どうか軍の端にでも<NL>お加えください<NL2><JUMP:31A8>"

b11_msg_002:
    "「あっ<PAGE><NAME_F6>様！<NL>劉璋様が涪水城で<NL>反乱軍に<PAGE>捕らわれてしまいました　どうか<NL>救い出してください<END>"

b11_msg_003:
    "「ここは涪水城です<END>"

b11_msg_004:
    "「この城のどこかに劉璋様が<NL>捕らわれているはずです<NL>助け出してください<END>"

b11_msg_005:
    "「雒城には蜀の猛将<NL>呉懿がいます<END>"

b11_msg_006:
    "「蜀にはものすごく硬い鉄が<NL>採れる山があるそうです<END>"

b11_msg_007:
    "「最後の虎が漢中軍に<NL>いるそうじゃ！<END>"

b11_msg_008:
    "「反乱軍が反乱しております<END>"

b11_msg_009:
    "「兄貴！この扉　体当たりで<NL>開けられそうだぜ！<END>"

b11_msg_010:
    "「<NAME_F6>殿　かたじけない<PAGE>一刻も早く反乱軍を<NL>鎮圧し　民を<NL>安んじてください<END>"

b11_msg_011:
    "「ここは雒城です<END>"

b11_msg_012:
    "「剛鉄山でございますか？<NL>それならここから南にまっすぐ<NL>行くとよいでしょう<END>"

b11_msg_013:
    "「昔は剛鉄山で採れる<NL>鉄からとてもよく切れる剣が<NL>作れたそうです<END>"

b11_msg_014:
    "「蜀には剣を作らせると右に<NL>出る者がいないほどの<NL>名刀鍛冶がいると<PAGE>聞くが--<END>"

b11_msg_015:
    "「あっしは剣なんぞ作れませんぜ！<END>"

b11_msg_016:
    "「剛鉄山を越えると<NL>綿竹関がございます<END>"

b11_msg_017:
    "「劉璋様は劉備軍を<NL>裏切り　漢中軍と手を<NL>組んだようです<PAGE>これほど世話になっておきながら<NL>情けない人だ<END>"

b11_msg_018:
    "「お前らも劉備軍に<NL>見つからないようにしろよ<END>"

b11_msg_019:
    "「わしは二十年も剣を<NL>作っておる<PAGE>しかし　いまだに納得のいく<NL>ものを作ったことがない<PAGE>おぬしに頼みがある<PAGE>剛鉄山に眠る鉄を探し<NL>わしに最強の剣を<NL>作らせてくれんじゃろうか？<END>"

b11_msg_020:
    "「おー　これじゃ　この輝き<NL>これこそわしの探し求めていた<NL>鉄じゃ<PAGE>うむ　早速剣を作ってみよう<END>"

b11_msg_021:
    "「少し時間がかかる　大事な<NL>用を済ませてから来て<NL>くれぬか？<END>"

b11_msg_022:
    "「す　すまぬ！！剣は<NL>出来上がったのだが<PAGE>黄巾賊の残党に盗まれて<NL>しまった！！！剣は全部で<NL>５本<PAGE>あの剣はおぬしらのものじゃ<PAGE>何とかして取り戻してくれ！<PAGE>剣の名は　屠龍の剣<NL>青龍刀　青釭の剣<NL>莫耶の剣<PAGE>干将の剣じゃ！<PAGE>すまぬ---<END>"

b11_msg_023:
    "「ここは綿竹関です<END>"

b11_msg_024:
    "「裏切り者の劉璋は<NL>成都城に立て籠もり<NL>劉備軍を待ち構えています<END>"

b11_msg_025:
    "「ぷれいやあを信じられない<NL>なんて---<END>"

b11_msg_026:
    "「うちの人　鉄を探しに<NL>剛鉄山へ行ったきり帰って<NL>こないのよ！３年も！<END>"

b11_msg_027:
    "「鉄なんかねーじゃねーか！<NL>かーちゃんに合わす顔がねーや<END>"

b11_msg_028:
    "「ここは成都城です<END>"

b11_msg_029:
    "「ここ　成都城は蜀の都<PAGE>どうぞ劉備様をお呼び<NL>ください<END>"

b11_msg_030:
    "「魏の曹操　呉の孫策<NL>蜀の劉備<PAGE>劉備殿も大きくなられたのー<NL>よしよし<END>"

b11_msg_031:
    "「むむ---不吉だ！<PAGE>帝星明らかならず　星座<NL>星漢　皆乱れる<PAGE>乱世はまだ続くぞ<END>"

b11_msg_032:
    "「曹操様がご病気になられ<NL>魏では大騒ぎだそうです<END>"

b11_msg_033:
    "<JUMP:2F0E>"

b11_msg_034:
    "「大至急荊州に<NL>向かってくれ！<END>"

b11_msg_035:
    "しかし　私の任された土地は<NL>広く　とても一人では目も<NL>届かぬ<PAGE>そこでどうだろう　荊州の<NL>長沙城を関羽が<PAGE>桂陽城を張飛が守っては<NL>くれまいか？<PAGE>今後は大きな戦も<NL>ないであろうし　二人とももう<NL>若くはない<PAGE>どうであろう　荊州を守っては<NL>くれぬか？<END>"

b11_msg_036:
    "「荊州は重要な土地<NL>おぬしらでなければ無事に<NL>治めることはできぬ<PAGE>二人で荊州に行って<NL>くれぬか？<END>"

b11_msg_037:
    "「荊州は蜀にとって重要な<NL>土地　二度と乱れぬよう<NL>良い政治を行ってくれ！<WAIT>関羽と張飛は荊州に<NL2>向かって旅立っていった<END>"

b11_msg_038:
    "「た　大変でございます　魏の<NL>曹操殿が病死いたしました<END>"

b11_msg_039:
    "「た　大変でございます　呉の<NL>孫策殿が暗殺されました<END>"

b11_msg_040:
    "「ぬぬ---なんと不吉な<NL>これではこの国は再び<NL>乱れるかもしれぬ<END>"

b11_msg_041:
    "「た　大変でございます　急に<NL>呉の軍勢が攻め込んで<NL>参りました<PAGE>関羽様も張飛様も<NL>ご自分の城を守るのが<NL>精一杯で<PAGE>呉の軍を追い返すことなど<NL>とてもできません<PAGE>何卒援軍をお出し<NL>くださいませ<END>"

b11_msg_042:
    "「関羽よ！張飛よ！よくぞ<NL>逆賊劉璋を討ってくれた<END>"

b11_msg_043:
    "「おお　<NAME_F6>　よくぞ<NL>来てくれた<PAGE>呉の軍に襲われ　この城も<NL>危ないところであった<PAGE>できれば一緒に戦いたいが<NL>荊州を任された身ではそうも<NL>いかぬ<PAGE>どうであろう　わしに代わって<NL>息子を連れていっては<NL>くれぬか？<END>"

b11_msg_044:
    "「そうだ　呉の内情について<NL>張飛が調べがついたと<NL>言っていたな<WAIT>「桂陽の城のことも気になる<NL>大至急行ってやってくれ<END>"

b11_msg_045:
    "「桂陽の城のことも気になる<NL>大至急行ってやってくれ<END>"

b11_msg_046:
    "「信じられん話だが　呉の<NL>孫権はどこからか玉璽を<NL>手に入れ　帝を名乗り<PAGE>魏の曹丕は漢朝の帝の<NL>位を剥奪し　自ら大魏の<NL>皇帝を名乗ったそうだ<PAGE>一度成都城に戻って<NL>劉備様と相談して参れ<END>"

b11_msg_047:
    "<JUMP:311E>"

b11_msg_048:
    "「<NAME_F7>と申します<NL>父に代わって漢王朝の<NL>ために戦いたいと思います<PAGE>どうか兵の端にでも加えて<NL>ください<NL2><JUMP:1CE9>"

b11_msg_049:
    "「ここは建安城です<END>"

b11_msg_050:
    "「ここから南には行けませんぞ<NL>鄱陽湖から水があふれ<NL>川ができたので渡れませんのじゃ<END>"

b11_msg_051:
    "「水があふれ出るところを大きな<NL>岩ででもせき止めれば<NL>いいのだが---<PAGE>そうすればまた南の町に<NL>行けるようになるんだがなー<END>"

b11_msg_052:
    "「孫策様が暗殺されたんです<END>"

b11_msg_053:
    "「水をせき止める？今日はその<NL>話をよく聞くなー<END>"

b11_msg_054:
    "「おぬし　水をせき止める<NL>方法を聞きたいか？<END>"

b11_msg_055:
    "「鄱陽山を特別な爆薬を<NL>使って爆破すれば<PAGE>その土砂できっと水を<NL>せき止められるはずじゃ<PAGE>爆薬は　青白く光る<NL>硝石と　１０００年以上<NL>生きている老木の枝<PAGE>それとわしの持っている硫黄を<NL>調合すればよい<PAGE>とりあえず老木の枝と<NL>硝石を持ってくれば　それで<NL>よい<END>"

b11_msg_056:
    "「おー　これじゃ　これさえあれば<NL>爆薬なんぞすぐにできる<PAGE>ちょちょいのちょい<PAGE>ほれ　できたぞ！<PAGE>しばし待たれい<END>"

b11_msg_057:
    "「ドッドーーーン<END>"

b11_msg_058:
    "突然ものすごい音が聞こえ<NL2>大地が揺れた　しかし別に<NL2>変化はない<END>"

b11_msg_059:
    "「これでよし　さっ　行きなされ<END>"

b11_msg_060:
    "「この城の南はついこの間まで<NL>草原だったんです<PAGE>しかし水軍提督<NL>周瑜殿が他国の侵略を<NL>防ぐために<PAGE>誰にも渡れぬ川を<NL>造られたのです<END>"

b11_msg_061:
    "「昔は北の洞窟で硝石が<NL>よく採れたものじゃ<END>"

b11_msg_062:
    "「へー！１０００年も生きている<NL>木をお探しかい！あんた<NL>見る目がないねー！<PAGE>うちの裏にある木　あれが<NL>そうじゃよ！<END>"

b11_msg_063:
    "「ここから南に向かうと<NL>会稽の砦があります<END>"

b11_msg_064:
    "「やっぱりわしは天才だ<END>"

b11_msg_065:
    "「ひゃー　干上がった---！<END>"

b11_msg_066:
    "「よし　よし<END>"

b11_msg_067:
    "「やっぱりありましたか<END>"

b11_msg_068:
    "「ここは会稽の砦です<END>"

b11_msg_069:
    "「呉城は孔明の兄<NL>諸葛瑾が守っております<END>"

b11_msg_070:
    "「えっ　建安城から<NL>来られたのですか？<PAGE>どうやって川を渡って<NL>来られたのですか？<END>"

b11_msg_071:
    "「ここは呉城です<END>"

b11_msg_072:
    "「鄱陽には砦がございます<END>"

b11_msg_073:
    "「鄱陽に向かうには山道を<NL>抜けなければなりません<END>"

b11_msg_074:
    "「確かにキラリと光ったような<NL>気がしたのですがーーー<END>"

b11_msg_075:
    "「山道でキラリと光るものを<NL>見たという村人が増えて<NL>いるんですよ<END>"

b11_msg_076:
    "「キラリと光る山道の<NL>光り物は　どうやら魔物が<NL>潜んでいるらしいのです<END>"

b11_msg_077:
    "「私の名は<NAME_F7><NL>呉の将軍である<PAGE>荊州に攻め入ったのも<NL>大義あってのこと　早々に<NL>ここより立ち去れ！<END>"

b11_msg_078:
    "「キラリと光るものがいる<NL>凍りつくほど冷たい光だ<PAGE>魔物が現れた<PAGE>ーーーーかと思うほどだった<NL>そこには１本の光り輝く剣が<NL>横たわっていた<PAGE>どうやらこの光は剣から<NL>発せられているようだ！<PAGE>張苞は剣を手に入れた<END>"

b11_msg_079:
    "「ここは鄱陽の砦です<END>"

b11_msg_080:
    "「ここから南東に向かうと<NL>新都城がございます<END>"

b11_msg_081:
    "「張苞よ　呂蒙じゃ　わしの<NL>言ったことが分からんと<NL>見えるな<PAGE>これ以上呉の国を<NL>侵すでない<END>"

b11_msg_082:
    "「先ほどまで孫権様も<NL>おられましたが　陸遜将軍を<NL>連れて建業に帰られました<END>"

b11_msg_083:
    "「ここは新都城です<END>"

b11_msg_084:
    "「これ以上呉の国を侵す<NL>ことは許さん　呂蒙の剣<NL>受けてみよ！<NL2><JUMP:3334>"

b11_msg_085:
    "「た---確かに強い！しかし<NL>太史慈にはまだまだ歯が<NL>立つまい<END>"

b11_msg_086:
    "「張苞将軍にお願いが<NL>ございます<PAGE>この城の南にある小屋に<NL>海賊が住み着いて<NL>しまったのです<PAGE>どうかこの海賊を退治して<NL>もらえないでしょうか<NL2><JUMP:32EC>"

b11_msg_087:
    "「退治していただいた暁には<NL>心ばかりのお礼をしとう<NL>ございます<PAGE>どうかよろしくお願いいたします<END>"

b11_msg_088:
    "「そうですか　残念です<END>"

b11_msg_089:
    "<JUMP:3301>"

b11_msg_090:
    "「ありがとうございました　村人も<NL>皆喜んでおります<END>"

b11_msg_091:
    "「三代続いた呉の国ももう<NL>終わりかもしれない<END>"

b11_msg_092:
    "「孫権様のお城ですか？<NL>それならこの城の北東に<NL>ございます<END>"

b11_msg_093:
    "「てめえら　何しに来やがった！<NL>百姓が集まったところで<NL>俺たちに勝てるわけがねえー！<PAGE>これでも食らって帰んな！<NL2><JUMP:3392>"

b11_msg_094:
    "「旦那あー　あっしらは百姓に<NL>戻って一からやり直します<END>"

b11_msg_095:
    "「す---すまねえ　もう悪さは<NL>しねえ　許してくれ！<END>"

b11_msg_096:
    "「<NAME_F6>よ！よくここまで<NL>来た！<PAGE>魏と我が軍は<NL>不可侵条約を結び<NL>共に蜀を滅ぼし<PAGE>天下の実権を握る<NL>はずであったが<PAGE>こうなれば討ち死にも覚悟<NL>せねばなるまい<PAGE>者ども　かかれ！<END>"

b11_msg_097:
    "「なかなかやるのお！若造よ---<NL>だがわしは他の者とは違うぞ<PAGE>それを今から教えてやる<END>"

b11_msg_098:
    "「兄　孫策を暗殺したのは<NL>わしだ---<PAGE>今思えば　皆魏の曹丕に<NL>だまされてやったこと<PAGE>兄者には悪いことをしたと<NL>思っている　兄者には<NL>冥土で--グフッ<WAIT>孫権は息を引き取った<WAIT>三代続いた呉の国の<NL2>歴史も今閉じようと<NL2>している----<NL2><JUMP:3385>"

b11_msg_099:
    "「蜀の民を安心させるためにも<NL>成都城に劉備様を<NL>お呼びください<END>"

b11_msg_100:
    "「ほー　そうか　気が向いたら<NL>聞きに来なされ<END>"

b11_msg_101:
    "<NAME_F6><WAIT>「我が君　ご決心を促したい<NL>ことがございます<PAGE>逆賊曹丕が帝の位を<NL>奪った今　劉備様自ら<NL>皇帝の位に就き<PAGE>祖廟の霊を慰め　またもって<NL>万民を安んずべきと考えます<PAGE>劉備<NL>「その話をするな！<PAGE>孔明<WAIT>「しかし漢の帝の血を引く<NL>我が君が進んで正統を<NL>継ぐべきです<PAGE>劉備<NL>「黙れ！！<PAGE>そちたちは余を末代まで<NL>不忠不義の人と<NL>呼ばせたいのか！<END>"

b11_msg_102:
    "<NAME_F6><NL2>「曹丕　孫権と我が君とは<NL2>同じではございませぬ<PAGE>漢朝を奪った逆賊<NL2>曹丕を誰が懲らしめます<PAGE>孔明<NL2>「このまま黙っていては　魏の<NL2>曹丕<PAGE>呉の孫権の即位を<NL2>正統なものとして認めたことに<NL2>なります<PAGE><NAME_F6><NL2>「今まで数々の逆賊が<NL2>皇帝の位を狙うことを<NL2>阻止せんと戦ってきたことが<PAGE>無意味となります<END>"

b11_msg_103:
    "劉備<NL2>「よく分かった！余は小さな<NL2>ことにこだわり過ぎたようじゃ！<PAGE>曹丕と孫権の即位は<NL2>認めぬと天下に知らしめるため<NL2>余は皇帝を名乗ろう！！<WAIT>かくして---建安２６年４月<WAIT>大礼台が築かれ<NL2>文武官の見守る中で<NL2>劉備は玉璽を受け<WAIT>ここに蜀の皇帝を天下に<NL2>宣した<END>"

b11_msg_104:
    "「おーい　俺だ　韓忠だ！早く<NL>橋を架けてくれ！<WAIT>「おう！今架けるぞ！<END>"

b11_msg_105:
    "<JUMP:23B6>"

b11_msg_106:
    "「劉備軍がすぐそこまで<NL>来ているそうだ　早く入れ<END>"

b11_msg_107:
    "「何者だ！-----<PAGE>韓忠の仲間か？！<END>"

b11_msg_108:
    "「何者だ　よそ者は帰れ<END>"

b11_msg_109:
    "「うぬが劉備軍の<NL><NAME_F6>か！<PAGE>よくも我が弟　袁術を---<PAGE>弟の仇だ　血祭りに<NL>上げろ！<END>"

b11_msg_110:
    "「わっはっははは---わしだよ<NL>袁術だ　この間の借りは<NL>返すぜ<END>"

b11_msg_111:
    "「<NAME_F6>将軍よ！よくぞ<NL>逆賊袁術を倒してくれた<PAGE>だが<NAME_F6>よ！<NL>安心してはならぬ<PAGE>袁術には兄がおる　きっとこの<NL>城に攻めてくるに違いない<PAGE>心してかかるがよい<NL2><JUMP:20EB><END>"

b11_msg_112:
    "「<NAME_F6>将軍よ　よくぞ<NL>蜀を平定してくれた　心から<NL>礼を言うぞ！<PAGE>ところでおぬしらは蜀の<NL>名刀師に会ったか？<PAGE>おぬしらの使っている剣では<NL>この後渡りきれまい<PAGE>一度会ってみてはどうだろう<NL>きっと良い剣を作って<NL>くれるはずじゃ<END>"

b11_msg_113:
    "「ここは建業です<END>"

b11_msg_114:
    "「ここから先は魏の国じゃ<NL>気を引き締めねば身を<NL>滅ぼすぞ！<END>"

b11_msg_115:
    "「孫権様って本当は良い<NL>人なのよ！<PAGE>でも曹丕様と交流を<NL>深めるにつれ　人が変わって<NL>いったのよ---<END>"

b11_msg_116:
    "「呉の国もこれで平和に<NL>なるでしょう<END>"

b11_msg_117:
    "「魏の広陵城はこの城の<NL>北東にございます<END>"

b11_msg_118:
    "「曹丕は各城砦に<NL>曹一族を置き<PAGE>劉備軍を<NL>待ち構えております<PAGE>気をつけてくださいませ<END>"

b11_msg_119:
    "「逆賊を討ち　国に平和を<NL>取り戻してください<END>"

b11_msg_120:
b11_msg_121:
    "「わしらの代わりに劉備様を<NL>助け　漢王朝復興の<NL>ためにも逆賊を討伐してくれ<END>"

b11_msg_122:
    "<JUMP:2D92>"

b11_msg_123:
    "「ついに荊州も一つに<NL>まとまりました<NL>ありがとうございました<END>"

b11_msg_124:
    "「やっと重い年貢から<NL>解放され　暮らしが楽に<NL>なります<END>"

b11_msg_125:
    "「劉備様　大変でございます<NL>実は---カクカクシカジカ<PAGE>なに！　呉の孫権　魏の<NL>曹丕ともに帝を名乗ったと<NL>いうのか<PAGE>漢朝の帝の位を<NL>剥奪したということは<PAGE>事実上漢王朝の<NL>崩壊を意味するでは<NL>ないか---<WAIT><JUMP:30AE>"

b11_msg_126:
    "「<NAME_F6>よ！　直ちに<NL>呉の孫権と魏の曹丕を<NL>討伐せよ！<END>"

b11_msg_127:
    "「お金が足りないあるよ<END>"

b11_msg_128:
    "「よし　出陣だ！　劉備軍を<NL>叩き潰せ！<END>"

b11_msg_129:
    "「うぬぬ　<NAME_F5>が<NL>相手してやる！<END>"

b11_msg_130:
    "「敵将はどこだ<NL><NAME_F5>が相手じゃ<END>"

b11_msg_131:
    "「それ　一気に叩き潰せ！<END>"

b11_msg_132:
    "「<NAME_F5>の恐ろしさを<NL>国に帰って劉備に伝えい！<END>"

b11_msg_133:
    "「今じゃ　敵が整わぬうちに<NL>蹴散らせ！<END>"

b11_msg_134:
    "「わしの名は<NAME_F5>　いざ<NL>勝負！<END>"

b11_msg_135:
    "「来たな　劉備軍　地獄の<NL>閻魔に会わせてやるぜ！<END>"

b11_msg_136:
    "「よくぞ帰ってきた！<PAGE>一刻も早く逆賊董卓を<NL>討ち　国に平和を<NL>戻してくれ<END>"

b11_msg_137:
    "「<NAME_F6>将軍<NL>無事で何よりだ<PAGE>帝の位を奪った袁術を<NL>のさばらせておくわけにはいかぬ<WAIT>「皆と力を合わせ　袁術を<NL>討て！<END>"

b11_msg_138:
    "「無事であったか！<NL><NAME_F6>将軍よ！<PAGE>この世は乱世の世<NL>争い事が絶えぬ<PAGE>帝の身を狙う袁紹を<NL>討つのもよし<PAGE>争い事の絶えぬ荊州を<NL>平定するのもよし<PAGE>それはおぬしらに任せるぞ！<END>"

b11_msg_139:
    "「おお　<NAME_F6>か　よく<NL>戻った<PAGE>たった今　均と名乗る<NL>少年が現れ<PAGE>兄が帰ったので至急来て<NL>くれるよう言っておったぞ！　さあ<NL>行くがよい　<NAME_F6>！<NL2><JUMP:2E95>"

b11_msg_140:
    "<JUMP:31EB>"

b11_msg_141:
    "「<NAME_F6>将軍　よくぞ<NL>無事で帰ってきた<PAGE>今　蜀は内と外より<NL>攻められ　明日をも知れぬ<NL>状態が続いておる<PAGE>私と劉璋とは同族の身<PAGE>一刻も早く反乱軍を<NL>鎮圧し　蜀に平和を<NL>取り戻してくれ<END>"

b11_msg_142:
    "「無事であったか！<NL><NAME_F6>将軍よ！<PAGE>呉の孫権が魏の曹丕と<NL>手を組み　荊州に手を<NL>伸ばしている<PAGE>これ以上民を苦しめては<NL>いかん<PAGE>まずは呉の孫権を討ってくれ<END>"

b11_msg_143:
    "「関興将軍よ！　魏の国は<NL>猛将が多いと聞く<PAGE>気を引き締めてかかれ！<END>"

b11_msg_144:
    "「漢王朝復興のためにも<NL>一刻も早く司馬懿を<NL>討ってくれ！！<END>"

b11_msg_145:
    "残念ですが<NL>歴史書<CALL:1000>は<NL>消えています<END>"

b11_msg_146:
    "どの歴史書を始めますか？<END>"

b11_msg_147:
    "文字の速さを決めてください<END>"

b11_msg_148:
    "どの歴史書を登録しますか？<END>"

b11_msg_149:
    "どの歴史書を消しますか？<END>"

b11_msg_150:
    "本当に宜しいですか？<END>"

b11_msg_151:
    "どの歴史書をこぴーしますか？<END>"

b11_msg_152:
    "どの歴史書にこぴーしますか？<END>"

b11_msg_153:
    "「魏の曹丕　呉の孫権<PAGE>またもや国は乱れる<END>"

b11_msg_154:
    "「エッ！　曹操様が<NL>亡くなられたって<END>"

b11_msg_155:
b11_msg_156:
b11_msg_157:
b11_msg_158:
b11_msg_159:
b11_msg_160:
    "「それもよかろう<END>"

b11_msg_161:
    "「よくやってくれたぞ<NL><NAME_F6>将軍よ！<PAGE>しかしいまだ荊州は<NL>平定されておらぬ<WAIT>「一刻も早く荊州を<NL>平定し　民を安んじてくれ<NL2><JUMP:20EB>"

b11_msg_162:
    "「一刻も早く荊州を<NL>平定し　民を安んじてくれ<NL2><JUMP:20EB>"

b11_msg_163:
    "「荊州平定　よくやってくれた<NL>民も安心するだろう<PAGE>しかし渤海の袁紹を<NL>討つのを忘れてはおらぬか<NL>「一刻も早く袁紹を<NL>討ってくれ！<END>"

b11_msg_164:
    "「一刻も早く袁紹を<NL>討ってくれ！<END>"

b11_msg_165:
    "「た　大変です<PAGE>趙雲は妹を冀州の<NL>袁紹にさらわれてしまい<PAGE>不本意ながら袁紹の<NL>武将として働いております<PAGE>妹を助け出し　趙雲を<NL>自由の身にしてやってください<END>"

b11_msg_166:
    "「そうですか---無事<NL>助け出せたのですネ<END>"

b11_msg_167:
    "「あー　どなたか存じませんが<NL>私をここから出してください<PAGE>でも---銀の鍵がないとこの<NL>扉は開きません<END>"

b11_msg_168:
    "「ありがとうございます<PAGE>広宗城にいる趙雲に<NL>会いましたら　この手紙を<NL>渡してください<WAIT><NAME_F6>は手紙を<NL2>受け取った<NL2><JUMP:28BC><END>"

b11_msg_169:
    "「ありがとうございます<END>"

b11_msg_170:
    "「<NAME_F6>将軍<NL>訳あって袁紹のもとで<NL>働いております<PAGE>将軍とは戦いたくありません<PAGE>どうかお引き取りください<END>"

b11_msg_171:
    "「こ　これは---<NL><NAME_F6>将軍<NL>かたじけない<PAGE>妹を助けていただいたのですね<PAGE>今度こそ将軍のために<NL>戦いたいと思います<PAGE>どうか兵の<NL>端にでも加えてください<NL2><JUMP:1CE9><END>"

b11_msg_172:
    "「将軍　鍵を買って<NL>ほしいのことあるよ　これ<NL>本当に役に立つよ<PAGE>１つ　金４０００でいいよ<NL2><JUMP:27A4>"

b11_msg_173:
    "「鍵　役立ちました？<END>"

b11_msg_174:
    "「得しますよ<END>"

b11_msg_175:
    "「損しますよ<END>"

b11_msg_176:
    "「<NAME_F6>将軍<NL>逆賊孫権をよくぞ<NL>倒してくれた<PAGE>しかし魏には帝の位を<NL>剥奪した張本人<PAGE>曹丕が軍備を蓄えたままで<NL>いる<PAGE>このままではこの国に真の<NL>平和は訪れぬ<PAGE>戦　戦の毎日で疲れては<NL>いようが　一刻も早く<NL>曹丕を倒し<PAGE>漢朝を復興させねばならぬ<PAGE>魏の国へは呉の建業より<NL>北東に攻め入るがよい<NL2><JUMP:20EB>"

b11_msg_177:
    "「ここは広陵城です<END>"

b11_msg_178:
    "「この中に諸葛亮様は<NL>おられますか？<NL2><JUMP:33B6>"

b11_msg_179:
    "「魏の<NAME_F7>という<NL>者です<PAGE>諸葛亮様は庵を出てより<NL>あまねく賢者を<NL>探し求めていると<PAGE>聞き及んでおります<PAGE>それがしに兵法のすべてを<NL>伝授していただきたく<NL>馳せ参じました<PAGE>どうか兵の<NL>端にでも加えてください<END>"

b11_msg_180:
    "「そうですか---残念です<END>"

b11_msg_181:
b11_msg_182:
    "「おらも戦に行くさ---<NL>ってなもんや三度笠---<NL>とくらあ！<END>"

b11_msg_183:
    "「うえーん　痛いよおー　あそこの<NL>木の下で足を<NL>切っちゃったよー！<END>"

b11_msg_184:
    "「あれっ　あれー　剣を<NL>なくしたみたいだぞー<PAGE>せっかく盗賊から<NL>巻き上げたのにもったいねえなあ<END>"

b11_msg_185:
    "「<NAME_F6>よ　劉備は<NL>蜀漢の皇帝を<NL>名乗ったそうじゃのお<PAGE>その昔　楼桑村で筵を<NL>売っていた若造が帝を<NL>名乗り<PAGE>さらに魏に攻め入るとは<NL>不届き千万<PAGE>まあ　いずれ蜀漢を叩く<NL>つもりであったからちょうどよいわ<PAGE>今この場で蜀軍の<NL>息の根を止めてくれる！<END>"

b11_msg_186:
    "「クッ　思ったよりやるわい　ここは<NL>一度退却して態勢を<NL>立て直せ！<PAGE>ダッダッダッダッ<END>"

b11_msg_187:
    "「ここは合肥の砦です<END>"

b11_msg_188:
    "「この砦の南西に汝南城が<NL>ございます<END>"

b11_msg_189:
    "「汝南城は元々<NL>曹彰将軍が<NL>守っておりました<PAGE>曹彰将軍は今は亡き<NL>曹操様の次男で<PAGE>武勇の優れた武将と<NL>聞いております<PAGE>決して侮ってはなりませんぞ<END>"

b11_msg_190:
    "「やはりここまで来たか　しかし<NL>わしもこの前とは違うぞ<PAGE>これでも食らえ！<END>"

b11_msg_191:
    "「たっ　大変です　司馬懿の<NL>軍が現れました<END>"

b11_msg_192:
    "「おー　司馬懿　よくぞ<NL>駆けつけてくれた　心より礼を<NL>言うぞ<PAGE><END>"

b11_msg_193:
    "「曹丕様---<PAGE>我らが劉備軍との戦いに<NL>一切参加しなかったわけを<NL>今明らかにいたしましょう<END>"

b11_msg_194:
    "司馬懿の攻撃<NL2>バシュッ<END>"

b11_msg_195:
    "「グオー　司馬懿<NL>貴様----<WAIT><END>"

b11_msg_196:
    "「うっ　裏切りかー<PAGE><END>"

b11_msg_197:
    "「いかん　このままでは全滅だ<NL>退却せよ<END>"

b11_msg_198:
    "「そうはいかん　とどめだ！<END>"

b11_msg_199:
    "「む---無念---<END>"

b11_msg_200:
    "「<NAME_F6>よ---曹丕の<NL>首はわしが取った<PAGE>しかし魏の国を貴様らにやる<NL>わけにはいかん　これを土産に<NL>建業まで帰るがいい<END>"

b11_msg_201:
    "「貴様らごときにこの<NL>司馬懿様の落雷をかわす<NL>ことはできまい<END>"

b11_msg_202:
    "「ウッ　ここはどこだ<PAGE>そうか　我々は司馬懿の<NL>クーデターを<NL>目の当たりにし----<PAGE>その後----落雷に<NL>打たれた<PAGE>しかしその後の記憶がない<PAGE>劉備様なら何か知っている<NL>はずだ<END>"

b11_msg_203:
    "「<NAME_F6>将軍　もう<NL>大丈夫か？<WAIT>「我が軍全滅の危機を<NL>救い　君らをここまで<NL>運んできたのは<PAGE>関羽将軍の次男<NL>関索将軍だ<PAGE>彼もまた漢朝の復興を<NL>心より願っておる<PAGE>礼がてら彼のいる<NL>長沙城まで行ってみては<NL>どうだ<PAGE>司馬懿の使う妖術の<NL>謎も解けるかもしれんぞ<NL2><JUMP:3432>"

b11_msg_204:
    "「我が軍全滅の危機を<NL>救い　君らをここまで<NL>運んできたのは<PAGE>関羽将軍の次男<NL>関索将軍だ<PAGE>彼もまた漢朝の復興を<NL>心より願っておる<PAGE>礼がてら彼のいる<NL>長沙城まで行ってみては<NL>どうだ<PAGE>司馬懿の使う妖術の<NL>謎も解けるかもしれんぞ<NL2><JUMP:3432>"

b11_msg_205:
    "「大切な話がある　孔明を<NL>隊列に加えてきてくれ！<END>"

b11_msg_206:
    "「<NAME_F6>将軍　もう<NL>元気になられましたか<PAGE><END>"

b11_msg_207:
    "「汝南城で司馬懿の<NL>使った妖術は<PAGE>孫子の書に記されている<NL>落雷の計であると思われます<PAGE>そうであればわたくしがその<NL>妖術封じの呪文を<NL>存じております<PAGE>上　上　下　下<NL>左　右　上　下<PAGE>司馬懿が呪文を<NL>唱え始めたらこちらもこの<NL>呪文を唱えることです<END>"

b11_msg_208:
    "「劉備軍よ　また現れたか<PAGE>しかし貴様らごときでこの<NL>司馬懿様の落雷の計を<NL>破ることはできまい<END>"

b11_msg_209:
    "「くそー　どうすれば落雷の計を<NL>破ることができるんだ<END>"

b11_msg_210:
    "「むっ　なぜ落雷が落ちんのだ<PAGE>こうなれば仕方がない<NL>力ずくで叩き潰してくれる！<END>"

b11_msg_211:
    "「フン！　貴様らがここまで<NL>切れるとは<NL>思ってもみんかったわい<PAGE>今日のところは落雷の計を<NL>破ったことを称えて　ひとまず<NL>退散するとするか<NL2><JUMP:343A>"

b11_msg_212:
    "「ここは汝南城です<END>"

b11_msg_213:
    "「北西に陳城という城が<NL>ありますが　道がないので<NL>誰も行けません<END>"

b11_msg_214:
    "「その昔　こんな話を<NL>聞いたことがあります<PAGE>陳城に通じる洞窟があり<NL>その洞窟の入り口はこの<NL>城にある---と-<END>"

b11_msg_215:
    "「うーむ　全くもって不可解だ<NL>洞窟を抜けても洞窟<NL>やめた---私には分からん<END>"

b11_msg_216:
    "「だっ　誰だ！　貴様は<NL><NAME_F6><PAGE>ふん　貴様らの中にも少しは<NL>切れる者がおるようだのう<PAGE>しかしここまでのこのこ<NL>やってきたことを後悔させてやる<END>"

b11_msg_217:
    "「むっ　無念---<END>"

b11_msg_218:
    "「ここは陳留の砦です<END>"

b11_msg_219:
    "「司馬懿には二人の<NL>息子がおりますが<PAGE>どちらも麒麟児---と<NL>呼ばれるほどの天才児だと<NL>聞いております<END>"

b11_msg_220:
    "「曹操様が亡くなられてから<NL>というもの　ありとあらゆる<NL>天災が続き<PAGE>さらに戦も相次ぎ　魏の<NL>民は今　食べるものにも<NL>事欠く有様です<END>"

b11_msg_221:
    "「あら？あなたたちって少しなまりが<NL>あるんじゃない　それってもしかして<NL>蜀のなまり？<PAGE>うふっ　カッワイー<END>"

b11_msg_222:
    "「貴様ら蜀軍か？そうか<NL>この城が攻められるようでは<PAGE>魏の国もそう長くは<NL>ないかもしれん<PAGE>しかし　この城を明け渡す<NL>わけにはいかん！<END>"

b11_msg_223:
    "「貴様らの勢いを止めることは<NL>できんのか---<NL2><JUMP:3442>"

b11_msg_224:
    "「ここは陳城　洛陽に最も<NL>近い城です<END>"

b11_msg_225:
    "「司馬懿は洛陽に<NL>１００万の兵とともに<NL>立てこもったと聞きます<END>"

b11_msg_226:
    "「追い詰められたとはいえ<NL>司馬懿は魏の<NL>天才軍師<PAGE>ここから洛陽までにはきっと<NL>兵が伏せてあるでしょう<END>"

b11_msg_227:
    "「その昔　洛陽は<NL>董卓将軍によって<NL>焼き払われましたが<PAGE>今では完全に復興され<NL>魏随一の城でございます<END>"

b11_msg_228:
    "「司馬懿の息子は<NL>司馬昭と司馬師と<NL>いいます<PAGE>この二人が策を弄せば<NL>１００万の官軍が逃げ場を<NL>失うと言われております<END>"

b11_msg_229:
    "「こっ　このままでは司馬懿に<NL>魏の国が乗っ取られてしまう<END>"

b11_msg_230:
    "「司馬懿を倒せば<NL>このゲームもおシバイだ<END>"

b11_msg_231:
    "「よくここまで来たな<NL><NAME_F6>将軍<PAGE>しかし　その悪運もここまでよ！<NL>ここでくたばっちまいな<END>"

b11_msg_232:
    "「ふんっ　なかなかやるのー　しかし<NL>この攻撃に耐えられるかな！<PAGE><END>"

b11_msg_233:
    "「なかなかしぶといのおー　そこまで<NL>命を懸けて　劉備が<NL>貴様らに何をくれる？<PAGE><END>"

b11_msg_234:
    "「我が父　司馬懿のもとで<NL>働けば　贅沢な暮らしが<NL>できるものを--<PAGE><END>"

b11_msg_235:
    "「これだけ言ってもわからんか<NL>今すぐ我が国に降れば<NL>重く用いようというものを--<PAGE><END>"

b11_msg_236:
    "「どうして我ら兄弟の攻撃に<NL>耐えることができたんだ<PAGE>我らにはわからん<END>"

b11_msg_237:
    "「<NAME_F6>よ！我が<NL>生涯の夢<NL>天下統一！！<PAGE>貴様らごときに邪魔は<NL>させん！！！<END>"

b11_msg_238:
    "「ソ---曹操様のもとで<NL>軍師として召し抱えられ<NL>曹丕様を裏切るまで<PAGE>長いようで短い<NL>人生だったような気がする<PAGE>クーデターは成功したはずなのに<PAGE>まさか貴様らごときに<NL>倒されるとは夢にも<NL>思わなかった<PAGE>し---しかしこれでこの国も<NL>一つにまとまったことになる<PAGE>こ---これからは---<NL>---ぐふっ---劉備と---<PAGE>ごふっ---<WAIT>司馬懿は息を引き取った<NL2><JUMP:345D>"

b11_msg_239:
    "「ここは洛陽　漢の都で<NL>ございます<END>"

b11_msg_240:
    "「司馬懿の軍はすべて倒され<NL>これからは平和が<NL>続くのでしょう<END>"

b11_msg_241:
    "「劉備様に復興した<NL>洛陽を見ていただきたい<NL>ものですね<END>"

b11_msg_242:
    "「あー---平和が訪れた<END>"

b11_msg_243:
    "「平和だと作物もよく<NL>とれるような気がするなあ<END>"

b11_msg_244:
    "「これでもう兵士として<NL>徴兵されることはないんだ--<NL>あー　青春だなあー<END>"

b11_msg_245:
    "「きっとパート２も出るんじゃ<NL>ないかなあー？<END>"

b11_msg_246:
    "「<CALL:346A>　<CALL:346E><NL><CALL:3472>　<CALL:3476><NL><CALL:347A>　<CALL:347E><PAGE><CALL:3482>　本当に<NL>ご苦労であった<PAGE>これでこの国がまた一つに<NL>まとまった<PAGE>もう二度と分裂することの<NL>ないよう　良い政治をする<NL>ことを誓おう<NL2><JUMP:348E>"

b11_msg_247:
b11_msg_248:
    "「関羽よ　俺は袁紹に<NL>仕えることにしたよ<END>"

b11_msg_249:
    "「きっ　貴様というやつは<END>"

b11_msg_250:
    "「くそ---悪運も尽きたか<NL>ぐふっ<WAIT>思えば　呂布の人生は<NL2>時代がそうさせたとはいえ<NL2>裏切りの人生であった<NL2><JUMP:31E3>"

b11_msg_251:
b11_msg_252:
b11_msg_253:
b11_msg_254:
b11_msg_255:
b11_msg_256:
    "「<NAME_F6>よ---この城<NL>他の城のようにはいかんぞ<END>"

b11_msg_257:
    "「むっ　無念<END>"

b11_msg_258:
b11_msg_259:
b11_msg_260:
b11_msg_261:
    "「龐統様ですが---たった今<NL>手勢を連れて出て<NL>いかれました<NL2><JUMP:2E81>"

b11_msg_262:
    "「荊州に攻め入るとは<NL>許しがたい　これでも食らえ！<END>"

b11_msg_263:
    "「なにっ　荊州を一つにまとめ<NL>平和を取り戻すために<NL>戦っておられたのか<PAGE>ならばどうか私を兵の端にでも<NL>加えてください<WAIT>龐統が我が軍に加わった<NL2><JUMP:2EA8>"

b11_msg_264:
    "「<NAME_F6>将軍<NL>大変です　<NAME_F5>が<NL>援軍に現れました<END>"

b11_msg_265:
    "「この洞窟を抜けたいなら<NL>身ぐるみ置いていきな<NL2><JUMP:28D7>"

b11_msg_266:
    "「--------<PAGE>--------<PAGE>--------<PAGE>ちょっと待った　貴殿の<NL>名前をお聞かせ願いたい<WAIT>「わしの名は<NAME_F6>！！<WAIT>「これはとんだ失礼をいたしました<WAIT>「私は<NAME_F7>という者<PAGE><NAME_F6>将軍に<NL>お仕えしたくて兵を集めて<NL>おりました<PAGE>どうか兵の端にでも<NL>加えてください<END>"

b11_msg_267:
    "「劉備軍の<NL><NAME_F6>将軍ですね<PAGE><NAME_F6>将軍がこの<NL>洞窟に入ったと聞き<NL>馳せ参じました<PAGE>どうか兵の端にでも<NL>加えてください<NL2><JUMP:1CE9>"

b11_msg_268:
    "小さな鍵穴を発見しました<END>"

b11_msg_269:
    "「私の名は<NAME_F7><NL>父に代わって逆賊と<NL>戦いとうございます<PAGE>どうか軍の端にでも<NL>お加えください<NL2><JUMP:1CE9><END>"

b11_msg_270:
    "「殿のお願いを聞き入れて<NL>くださいませ<END>"

b11_msg_271:
    "「殿は奥の部屋に---<END>"

b11_msg_272:
    "これで荊州にも平和が<NL2>訪れました<NL2><JUMP:2E62>"

b11_msg_273:
    "「<NAME_F6>よ<NL>ご苦労であった<PAGE>おー　そうじゃ　徐庶の<NL>母上の手紙を預かっておる<WAIT>徐庶は手紙を受け取った<WAIT>「--------<PAGE>--------<PAGE>--------<PAGE>--------<WAIT>「身寄りのない都で寂しい<NL>思いをしているのでしょう<PAGE>お願いでございます　私に<NL>お暇をくださいませ<WAIT>「徐庶よ　行って母上を<NL>安心させてあげなさい<WAIT>「あっ-ありがとうございます<END>"

b11_msg_274:
b11_msg_275:
    "<JUMP:2C1F>"

b11_msg_276:
    "「孔明の庵です　どうやら<NL>留守のようですなー<END>"

b11_msg_277:
    "「<NAME_F6>将軍！よくぞ<NL>帰ってきた！<PAGE>孔明がお前たちに話が<NL>あるそうだ　聞いてやってくれ！<END>"

b11_msg_278:
    "「聞いてのとおりだ！劉璋殿は<NL>私と同族である<PAGE><NAME_F6>よ　直ちに<NL>荒れ果てた蜀を救いに<NL>行ってくれ！<END>"

b11_msg_279:
    "劉備軍により　蜀も無事<NL2>平定された<NL2><JUMP:2F06>"

b11_msg_280:
    "<JUMP:2D09>"

b11_msg_281:
    "「ここは張飛様のお城です<NL>どなたもお呼びすることは<NL>できません<END>"

b11_msg_282:
    "「劉備に会って伝えろ　呉の<NL>後ろには魏がついているとな！<NL2><JUMP:318A>"

b11_msg_283:
    "「この城を取り戻したからと<NL>いって　３代続いた我らの<NL>呉が敗れたわけではない<PAGE>漢王朝の復興など<NL>させるものか！<NL2><JUMP:3169>"

b11_msg_284:
    "<JUMP:2EDE>"

b11_msg_285:
    "「<NAME_F6>将軍！<NL>涪水城に劉璋様が<NL>捕らわれていると聞きます<PAGE>一刻も早く助け出して<NL>くださいませ<END>"

b11_msg_286:
    "「ここから先は劉璋様の国<NL>蜀でございます<PAGE>争いごとのない　とても<NL>住みやすい国でございます<END>"

b11_msg_287:
    "「これは<NAME_F6>将軍<PAGE>張松が孔明殿に申した<NL>とおり　蜀の国は乱れきって<NL>おります<PAGE>どうか反乱軍を鎮圧し<NL>蜀の国に再び平和を<NL>取り戻してください<END>"

b11_msg_288:
    "劉備軍の防御力が<NL2>上がった<END>"

b11_msg_289:
    "<NAME_F5>軍の防御力が<NL2>上がった<END>"

b11_msg_290:
b11_msg_291:
    "火計に対する防御力を<NL2>高めた<END>"

b11_msg_292:
b11_msg_293:
    "水計に対する防御力を<NL2>高めた<END>"

b11_msg_294:
b11_msg_295:
    "敵の使用する策略を<NL2>無効にした<END>"

b11_msg_296:
b11_msg_297:
    "敵の攻撃を一切<NL2>受け付けなくなった<END>"

b11_msg_298:
    "城のランクが下がった<END>"

b11_msg_299:
    "城のランクが上がった<END>"

b11_msg_300:
    "敵将<NAME_F4>に<NL2>疑心暗鬼を起こさせた<END>"

b11_msg_301:
    "<NAME_F4>は疑心暗鬼を<NL2>起こした<END>"

b11_msg_302:
    "敵将<NAME_F4>に離間を<NL2>起こさせた<END>"

b11_msg_303:
    "<NAME_F4>は離間を<NL2>起こした<END>"

b11_msg_304:
b11_msg_305:
    "<NAME_F4>の素早さが<NL2>上がった<END>"

b11_msg_306:
    "<NAME_F4>軍の使った<NL2>兵力すべて解き放った<END>"

b11_msg_307:
    "劉備軍の使った兵力すべて<NL2>解き放った<END>"

b11_msg_308:
    "<NAME_F5>は<NAME_F4>の<NL2>首を取った<END>"

b11_msg_309:
    "<NAME_F4>は<NAME_F5>に<NL2>首を取られた<END>"

b11_msg_310:
b11_msg_311:
    "<NAME_F4>の攻撃力が<NL2>上がった<END>"

b11_msg_312:
b11_msg_313:
    "<NAME_F5>に伏兵がついた<END>"

b11_msg_314:
    "劉備軍は退却した<END>"

b11_msg_315:
    "<NAME_F5>軍は退却した<END>"

b11_msg_316:
    "劉備軍は退却した<END>"

b11_msg_317:
    "<NAME_F5>軍は退却した<END>"

b11_msg_318:
    "伏兵の攻撃<END>"

b11_msg_319:
    "<NAME_F5>は<CALL:D565=さくりゃく>を<NL2>解き放った<END>"

b11_msg_320:
    "<NAME_F5>軍の<NL2>策免の計の効力が切れた<END>"

b11_msg_321:
    "劉備軍の策免の計の<NL2>効力が切れた<END>"

b11_msg_322:
    "<NAME_F5>軍の<NL2>撃免の計の効力が切れた<END>"

b11_msg_323:
    "劉備軍の撃免の計の<NL2>効力が切れた<END>"

b11_msg_324:
b11_msg_325:
b11_msg_326:
b11_msg_327:
b11_msg_328:
    "これで荊州と蜀にも平和な<NL>日々が戻るはずだ<PAGE>この国も　北の魏は<NL>曹操殿が帝を補佐し<NL>東の呉は孫策殿が<PAGE>そして今　西の蜀を私が<NL>治めることとなった<PAGE>これからも漢王朝のため<NL>帝のため　そして私を<NL>信じてくれた蜀の民のためにも<PAGE>我らは戦い続けねばならぬ！<END>"

b11_msg_329:
    "「海賊どもを退治していただき<NL>ありがとうございます<PAGE>お礼の品を用意いたしました<PAGE>これは我々が海賊の<NL>一人を捕らえたときに<NL>手に入れた青龍刀です<PAGE>ご覧ください　この輝きはただの<NL>青龍刀ではありません<PAGE>正義のために使ってください<END>"

b11_msg_330:
    "「聞いてのとおりだ　関羽も<NL>張飛も私の義兄弟<PAGE>大至急荊州に<NL>向かってくれ！<END>"

b11_msg_331:
    "<JUMP:2D3E>"

b11_msg_332:
b11_msg_333:
b11_msg_334:
b11_msg_335:
b11_msg_336:
b11_msg_337:
b11_msg_338:
b11_msg_339:
b11_msg_340:
b11_msg_341:
b11_msg_342:
b11_msg_343:
b11_msg_344:
b11_msg_345:
b11_msg_346:
b11_msg_347:
b11_msg_348:
b11_msg_349:
b11_msg_350:
b11_msg_351:
b11_msg_352:
b11_msg_353:
b11_msg_354:
b11_msg_355:
b11_msg_356:
b11_msg_357:
b11_msg_358:
b11_msg_359:
b11_msg_360:
b11_msg_361:
b11_msg_362:
b11_msg_363:
b11_msg_364:
b11_msg_365:
b11_msg_366:
b11_msg_367:
b11_msg_368:
b11_msg_369:
b11_msg_370:
b11_msg_371:
b11_msg_372:
b11_msg_373:
b11_msg_374:
b11_msg_375:
b11_msg_376:
b11_msg_377:
b11_msg_378:
b11_msg_379:
b11_msg_380:
b11_msg_381:
b11_msg_382:
b11_msg_383:
    ""
